% Chapter I: Basic Concepts
% Transcribed from PDF pages 17--24 (book pages 1--8)

\chapter{Basic Concepts}
\label{ch:1}

\section{Introduction}
\label{sec:1-1}

The equations of hydrodynamics, in spite of their complexity, allow
some simple patterns of flow (such as between parallel planes, or rotating
cylinders) as stationary solutions.  These patterns of flow can, however,
be realized only for certain ranges of the parameters characterizing them.
Outside these ranges, they cannot be realized.  The reason for this lies
in their inherent \emph{instability}, i.e.\ in their inability to sustain themselves
against small perturbations to which any physical system is subject.  It
is in the differentiation of the stable from the unstable patterns of
permissible flows that the problems of \emph{hydrodynamic stability} originate.

In recent years the class of such problems of stability has been enlarged
by the interest in hydrodynamic flows of electrically conducting fluids
in the presence of magnetic fields.  This is the domain of \emph{hydromagnetics};
and there are problems of \emph{hydromagnetic stability} even as there are
problems of hydrodynamic stability.

This book is devoted to a consideration of some typical problems in
hydrodynamic and hydromagnetic stability; and in this introductory
chapter we shall formulate the basic principles and concepts of the
subject.

\section{Basic concepts}
\label{sec:1-2}

Suppose, then, that we have a hydrodynamic system which, in
accordance with the equations governing it, is in a \emph{stationary} state, i.e.\
in a state in which none of the variables describing it is a function of
time.  Let $X_1, X_2, \ldots, X_j$ be a set of parameters which define the system.
These parameters will include geometrical parameters such as the
dimensions of the system; parameters characterizing the velocity field
which may prevail in the system; the magnitudes of the forces which
may be acting on the system, such as pressure gradients, temperature
gradients, magnetic fields and rotation; and others.

In considering the stability of such a system (with a given set of
parameters $X_1, \ldots, X_j$) we essentially seek to determine the reaction of
the system to small disturbances.  Specifically, we ask: if the system is
disturbed, will the disturbance gradually die down, or will the disturbance
grow in amplitude in such a way that the system progressively departs
from the initial state and never reverts to it?  In the former case, we say
that the system is \emph{stable} with respect to the particular disturbance; and
in the latter case, we say that it is \emph{unstable}.  Clearly, a system must be
considered as unstable even if there is only \emph{one} special mode of disturbance with respect to which it is unstable.  And a system cannot be
considered stable unless it is stable with respect to \emph{every} possible disturbance to which it can be subject.  In other words, stability must imply
that there exists no mode of disturbance for which it is unstable.

If all initial states are classified as stable, or unstable, according to
the criteria stated, then in the space of the parameters, $X_1, \ldots, X_j$, the
locus which separates the two classes of states defines the states of
\emph{marginal stability} of the system.  By this definition, a \emph{marginal state} is
a state of \emph{neutral stability}.

The locus of the marginal states in the $(X_1, \ldots, X_j)$-space will be defined
by an equation of the form
\begin{equation}
\label{eq:1-1}
\Sigma(X_1, \ldots, X_j) = 0.
\end{equation}
The determination of this locus is one of the prime objects of an investigation on hydrodynamic stability.

In thinking of the stability of a hydrodynamic system, it is often
convenient to suppose that all parameters of the system, save one, are
kept constant while the chosen one is continuously varied.  We shall
then pass from stable to unstable states when the particular parameter
we are varying takes a certain critical value.  We may say that \emph{instability}
sets in at this value of the chosen parameter when all the others have
their preassigned values.

States of marginal stability can be of one of two kinds.  The two kinds
correspond to the two ways in which the amplitudes of a small disturbance can grow or be damped: they can grow (or be damped) aperiodically; or they can grow (or be damped) by oscillations of increasing (or
decreasing) amplitude.  In the former case, the transition from stability
to instability takes place via a marginal state exhibiting a stationary
pattern of motions.  In the latter case, the transition takes place via a
marginal state exhibiting oscillatory motions with a certain definite
characteristic frequency.

If at the onset of instability a stationary pattern of motions prevails,
then one says that the \emph{principle of the exchange of stabilities} is valid and
that instability sets in as stationary \emph{cellular convection}, or \emph{secondary
flow}.  (The significance of the qualification `cellular' will become clear
presently.)  On the other hand, if at the onset of instability oscillatory
motions prevail, then one says (following Eddington) that one has a
case of \emph{overstability}.  Eddington explains this choice of the terminology
as follows: `In the usual kinds of instability, a slight displacement
provokes restoring forces tending away from equilibrium; in overstability
it provokes restoring forces so strong as to overshoot the corresponding
position on the other side of equilibrium.'

In classifying marginal states into two classes---stationary and
oscillatory---we have supposed that we are dealing with \emph{dissipative
systems}.  In non-dissipative, conservative, systems the situation is
generally somewhat different.  In these cases the stable states, when
perturbed, execute undamped oscillations with certain definite characteristic frequencies; while in the unstable states small initial perturbations tend to grow exponentially with time; and the marginal states
themselves are stationary.

\section{The analysis in terms of normal modes}
\label{sec:1-3}

The mathematical treatment of a problem in stability generally
proceeds along the following lines.

We start from an initial flow which represents a stationary state of
the system.  By supposing that the various physical variables describing
the flow suffer small (infinitesimal) increments, we first obtain the
equations governing these increments.  In obtaining these equations
from the relevant equations of motion, we neglect all products and
powers (higher than the first) of the increments and retain only terms
which are linear in them.  The theory derived on the basis of such
\emph{linearized equations} is called the linear stability theory in contrast to
\emph{non-linear theories} which attempt to allow for the finite amplitudes of
the perturbations.  In this book we shall be concerned only with the
linear stability theory which presupposes that the perturbations are
infinitesimal; and this must be understood even if the qualification is
not explicitly made every time.

As we have explained, stability means stability with respect to all
possible (infinitesimal) disturbances.  Accordingly, for an investigation
on stability to be complete, it is necessary that the reaction of the system
to all possible disturbances be examined.  In practice, this is accomplished
by expressing an arbitrary disturbance as a superposition of certain
basic possible modes and examining the stability of the system with
respect to each of these modes.  In illustration of how one does this,
consider a system confined between two parallel planes and in which
the physical variables in the stationary state are functions only of the
coordinate ($z$, say) normal to the planes.  In this case, we may analyse
an arbitrary disturbance in terms of two-dimensional periodic waves.
Thus, if $A(x, y, z, t)$ represents a typical amplitude describing the
disturbance, we expand it in the manner
\begin{equation}
\label{eq:1-2}
A(x, y, z, t) = \int_{-\infty}^{+\infty} \int_{-\infty}^{+\infty}
dk_x\, dk_y\, A_k(z, t) \exp[i(k_x x + k_y y)],
\end{equation}
where
\begin{equation}
\label{eq:1-3}
k = \sqrt{k_x^2 + k_y^2}
\end{equation}
is the wave number associated with the disturbance $A_k(z, t)$.  Since the
perturbation equations are linear, the reaction of the system to a general
disturbance can be determined if we know the reaction of the system to
disturbances of all assigned wave numbers.  In particular, the stability
of the system will depend on its stability to disturbances of \emph{all} wave
numbers; and instability will follow from its instability to disturbances
of even one wave number.  From this discussion it is clear that for `plane
problems' of the kind considered, the marginal state will be of neutral
stability for disturbances of one particular wave number ($k_c$, say) and
will be stable for disturbances of all other wave numbers.  At the onset
of instability, therefore, motions (stationary or oscillatory, as the case
may be) will be manifested which will belong to a particular wave
number (namely, $k_c$); they will accordingly present a cellular pattern,
the horizontal dimension of the cell being determined by the wavelength
$\lambda_c = 2\pi/k_c$.

In problems with other geometries, the disturbance will have to be
analysed correspondingly.  Thus, if the stationary state has symmetry
about an axis, we may expand the disturbance in the manner
\begin{equation}
\label{eq:1-4}
A(\varpi, z, \theta, t) = \sum_{m=-\infty}^{+\infty} \int_{-\infty}^{+\infty}
A_{m,k}(\varpi, t) e^{i(kz + m\theta)}\, dk
\end{equation}
(where $z$ is measured along the axis of symmetry, $\varpi$ is the distance from
the axis, and $\theta$ is the azimuthal angle), and investigate the stability of the
system with respect to all modes distinguished by $k$ and $m$.  Similarly,
if the initial state has spherical symmetry, we may expand the disturbance in spherical harmonics, $Y_l^m(\theta, \varphi)$, in the manner
\begin{equation}
\label{eq:1-5}
A(r, \theta, \varphi) = \sum_{l=0}^{\infty} \sum_{m=-l}^{+l}
Y_l^m(\theta, \varphi) A_l^m(r, t)
\end{equation}
(where $r$ denotes the distance from the origin, and $\theta$ and $\varphi$ are the
spherical polar angles), and investigate the stability of the system with
respect to all modes distinguished by $m$ and $l$.

The essential point of the foregoing remarks is that in all cases the
disturbance must be expanded in terms of some suitable set of \emph{normal
modes} which must be \emph{complete} for such an expansion to be possible.  Let
the various modes appropriate to a particular problem be distinguished
by the symbol~$\mathbf{k}$.  In practice, several parameters may be needed to
distinguish the different modes; and the symbol~$\mathbf{k}$ is assumed to represent
all the parameters that may be needed.  With this understanding, we
may write (again, symbolically)
\begin{equation}
\label{eq:1-6}
A(\mathbf{r}, t) = \int A_{\mathbf{k}}(\mathbf{r}, t)\, d\mathbf{k}.
\end{equation}

The equations governing the general (infinitesimal) perturbations
can be specialized for the normal modes.  We then eliminate the dependence on time by seeking solutions of the form
\begin{equation}
\label{eq:1-7}
A_{\mathbf{k}}(\mathbf{r}, t) = A_{\mathbf{k}}(\mathbf{r}) e^{p_{\mathbf{k}} t},
\end{equation}
where $p_{\mathbf{k}}$ is a constant to be determined.  The subscript~$\mathbf{k}$ has been
attached to $p$ to emphasize that the value of this constant will be different
for the different normal modes distinguished by~$\mathbf{k}$.

With the dependence on time separated in this manner, the perturbation equations will involve $p_{\mathbf{k}}$ as a parameter.  And solutions of these last
equations must be sought which satisfy certain necessary boundary
conditions (such as, that there be no slip at a `rigid boundary', or that
there be no tangential viscous stresses at a `free boundary').  In general
the equations will not allow non-trivial solutions (i.e.\ not vanishing
everywhere) for any arbitrarily assigned~$p_{\mathbf{k}}$.  Indeed, the requirement
that the equations allow non-trivial solutions satisfying the various
boundary conditions leads directly to a \emph{characteristic value problem} for
$p_{\mathbf{k}}$.  The problem is thus reduced to determining $p_{\mathbf{k}}$ for the various modes.
In general the characteristic values for $p_{\mathbf{k}}$ will be complex:
\begin{equation}
\label{eq:1-8}
p_{\mathbf{k}} = p_{\mathbf{k}}^{(r)} + i p_{\mathbf{k}}^{(i)},
\end{equation}
where $p_{\mathbf{k}}^{(r)}$ and $p_{\mathbf{k}}^{(i)}$ will depend, apart from~$\mathbf{k}$, on the parameters $X_1, \ldots, X_j$
of the basic flow.  The condition for stability is that $p_{\mathbf{k}}^{(r)}$ be negative for
all~$\mathbf{k}$.  The states which will be of neutral stability with respect to
disturbances belonging to~$\mathbf{k}$ will be characterized by
\begin{equation}
\label{eq:1-9}
p_{\mathbf{k}}^{(r)}(X_1, \ldots, X_j) = 0.
\end{equation}
This is a condition on the parameters $X_1, \ldots, X_j$; and it defines a locus,
\begin{equation}
\label{eq:1-10}
\Sigma_{\mathbf{k}}(X_1, \ldots, X_j) = 0,
\end{equation}
in the $(X_1, \ldots, X_j)$-space; this locus separates states which are stable
from those which are unstable with respect to disturbances belonging
to the particular mode~$\mathbf{k}$.  It follows from this that the \emph{locus $\Sigma(X_1, \ldots, X_j)$
separating the regions of complete stability from the regions of instability
in the $(X_1, \ldots, X_j)$-space is the envelope of the loci $\Sigma_{\mathbf{k}}$}; further, \emph{that when the
system becomes unstable as we cross $\Sigma$ at some particular point, the mode
of disturbance which will be manifested (at onset) will be the one whose
locus $\Sigma_{\mathbf{k}}$ touches $\Sigma$ at that point.}

A further observation might be made.  The distinction between the
two kinds of marginal states (stationary and oscillatory) made in \S\,2
corresponds to whether or not the imaginary part, $p_{\mathbf{k}}^{(i)}$ of $p_{\mathbf{k}}$, vanishes
when the real part, $p_{\mathbf{k}}^{(r)}$ of $p_{\mathbf{k}}$, does.  If $p_{\mathbf{k}}^{(i)} = 0$ implies that $p_{\mathbf{k}}^{(r)} = 0$ for
every~$\mathbf{k}$, then the principle of the exchange of stabilities will be valid;
otherwise, we will have overstability at least when instability sets in as
certain modes.  If overstability occurs, then $p_{\mathbf{k}}^{(i)}$ will have determinate
values on~$\Sigma_{\mathbf{k}}$.  Accordingly, in case of overstability, the theory must tell
us not only what the mode of disturbance is that will be manifested at
onset of instability, but also the characteristic frequency of the oscillations.

\section{Non-dimensional numbers}
\label{sec:1-4}

As we have explained, the solution of a problem in stability generally
centres on the determination of the locus~$\Sigma$ of the marginal states in the
space of the parameters $\{X_1, \ldots, X_j\}$ characterizing the initial flow.  In
expressing the results of the theory, it will often be convenient to combine
the various parameters into certain non-dimensional combinations or
\emph{numbers}.  The most famous of such numbers is the \emph{Reynolds number}~$\mathrm{R}$
which is a combination of the form
\begin{equation}
\label{eq:1-11}
\mathrm{R} = Lv/\nu,
\end{equation}
where $L$ is a typical dimension of the system, $v$ is a measure of the
velocities which prevail in the stationary flow, and $\nu$ is the kinematic
viscosity.  However, particular problems may require the definition of
special numbers.  Thus, in the problem of the stability of viscous flow
between two rotating coaxial cylinders, the parameters are the radii
$R_1$ and $R_2$ of the two cylinders, the angular velocities $\Omega_1$ and $\Omega_2$ of their
rotations, and the kinematic viscosity~$\nu$.  The non-dimensional combinations of these parameters in terms of which it is convenient to discuss
this problem are (see Chapter~VII):
\begin{equation}
\label{eq:1-12}
\eta = \frac{R_1}{R_2},\quad \mu = \frac{\Omega_2}{\Omega_1},
\quad\text{and}\quad
\mathscr{T} = \frac{4\Omega_1^2 R_1^4}{\nu^2}
\cdot \frac{(1-\mu)(1-\mu/\eta^2)}{(1-\eta^2)^2}.
\end{equation}

The combination
\begin{equation}
\label{eq:1-13}
\frac{4\Omega^2 L^4}{\nu^2}
\end{equation}
of an angular velocity~$\Omega$, a linear dimension~$L$, and the kinematic
viscosity~$\nu$, which occurs in the definition of~$\mathscr{T}$, always appears when we
are dealing with rotating systems; it is called the \emph{Taylor number}.
Similarly, in the problem of the thermal stability of a horizontal layer
of fluid heated from below, the combinations which occur are:
\begin{equation}
\label{eq:1-14}
R = \frac{g\alpha\beta}{\kappa\nu}d^4 \quad\text{and}\quad p = \nu/\kappa,
\end{equation}
where $g$ is the acceleration due to gravity, $d$ is the depth of the layer,
$\beta$ is the adverse temperature gradient which is maintained, and
$\alpha$, $\kappa$, and $\nu$ are the coefficients of volume expansion, thermometric
conductivity, and kinematic viscosity, respectively; $R$ as defined here
is called the \emph{Rayleigh number}; and $p$ is the \emph{Prandtl number}.  And when
we are dealing with an electrically conducting fluid in which a magnetic
field $H$ prevails, a `number' of frequent occurrence is
\begin{equation}
\label{eq:1-15}
Q = \frac{\mu_e H^2 \sigma}{\rho\nu}d^2,
\end{equation}
where $\mu_e$ denotes the magnetic permeability, $\rho$ is the density, $\sigma$ is the
coefficient of electrical conductivity, and $d$ is a linear dimension.

It should be pointed out that there is nothing unique in the way these
various numbers are defined; they happen to be the ones which have been
chosen; and in some sense they occur most naturally in certain types of
problems.

\section*{Bibliographical Notes}
\label{sec:1-bib}

The only book which is devoted exclusively to problems in hydrodynamic and
hydromagnetic stability is:

\begin{enumerate}
\item C.~C.~\textsc{Lin}, \emph{The Theory of Hydrodynamic Stability}, Cambridge, England, 1955.

A major portion of Lin's book treats problems which are related to viscous shear
flows and the Orr--Sommerfeld equation: topics which on account of their very
special character will not be considered in this book.
\end{enumerate}

\noindent
Problems in hydromagnetic stability are also considered in:

\begin{enumerate}\setcounter{enumi}{1}
\item T.~G.~\textsc{Cowling}, \emph{Magnetohydrodynamics}, Interscience Tracts on Physics
  and Astronomy, No.~4, Interscience Publishers, Inc., New York, 1957;
  see particularly chapter~4.
\end{enumerate}

\noindent
Among other books on related subjects, we may refer to:

\begin{enumerate}\setcounter{enumi}{2}
\item H.~\textsc{Alfv\'en}, \emph{Cosmical Electrodynamics}, International Series of Monographs
  on Physics, Oxford, England, 1950.
\item \textsc{Lyman Spitzer, Jr.}, \emph{Physics of Fully Ionized Gases}, Interscience Tracts
  on Physics and Astronomy, No.~3, Interscience Publishers, Inc., New York,
  1956.
\item J.~W.~\textsc{Dungey}, \emph{Cosmic Electrodynamics}, Cambridge Monographs on
  Mechanics and Applied Mathematics, Cambridge, England, 1958.
\end{enumerate}

\noindent
The following review articles bear on topics treated in this book:

\begin{enumerate}\setcounter{enumi}{5}
\item S.~\textsc{Lundquist}, `Studies in magneto-hydrodynamics', \emph{Arkiv f\"or Fysik}, \textbf{5},
  297--347 (1952).
\item T.~G.~\textsc{Cowling}, `Solar electrodynamics', \emph{The Sun}, chapter~8, edited by
  G.~P.~Kuiper, University of Chicago Press, Chicago, 1953.
\item W.~M.~\textsc{Elsasser}, `Hydromagnetism.\ I.\ A review', \emph{American J.\ Physics},
  \textbf{23}, 590--609 (1955); `Hydromagnetism.\ II.\ A review', ibid.\ \textbf{24}, 85--110
  (1956).
\item G.~H.~A.~\textsc{Cole}, `Some aspects of magnetohydrodynamics', \emph{Advances in
  Physics}, \textbf{5}, 452--497 (1956).
\end{enumerate}

\noindent
This book represents an expansion of the point of view taken in:

\begin{enumerate}\setcounter{enumi}{9}
\item S.~\textsc{Chandrasekhar}, `Problems of stability in hydrodynamics and hydromagnetics', \emph{Monthly Notices Roy.\ Astron.\ Soc.\ London}, \textbf{113}, 667--78 (1953).
\item S.~\textsc{Chandrasekhar}, `Thermal convection', \emph{Daedalus}, \textbf{86}, 323--39 (1957).
\end{enumerate}

\noindent
\S\,2.\ The quotation from Eddington is from:

\begin{enumerate}\setcounter{enumi}{11}
\item A.~S.~\textsc{Eddington}, \emph{The Internal Constitution of the Stars}, p.~201, Cambridge,
  England, 1926.
\end{enumerate}
